\documentclass[../../main.tex]{subfiles}

\begin{document}

\chapter{Register 8086 dan Register Flag}
\label{chap:register}

% =========================================================
% OUTCOME PEMBELAJARAN KHUSUS BAB
% =========================================================

\begin{subcpmkbox}
\SubCPMK{1.2} --- Mengidentifikasi fungsi register umum (\reg{AX},
\reg{BX}, \reg{CX}, \reg{DX}), register segmen (\reg{CS}, \reg{DS},
\reg{SS}, \reg{ES}), register \eng{pointer}/\eng{index} (\reg{SP},
\reg{BP}, \reg{SI}, \reg{DI}), dan register \eng{flag} beserta
pengaruhnya terhadap eksekusi program.

\vspace{0.4em}
\textbf{CPMK yang didukung:} \CPMK{01} \quad\textbullet\quad
\textbf{CPL yang didukung:} \CPL{02}
\end{subcpmkbox}

\begin{tujuanbox}
Setelah mempelajari bab ini, mahasiswa mampu:
\begin{enumerate}
    \item menyebutkan register umum 8086 beserta akses byte rendah
          dan byte tingginya;
    \item menjelaskan peran khas setiap register umum dan alasan
          historis di balik peran tersebut;
    \item menjelaskan fungsi register segmen dalam pengalamatan
          memori;
    \item menjelaskan peran register \eng{pointer} dan \eng{index}
          serta register \reg{IP};
    \item mengidentifikasi sembilan \eng{flag} pada register
          \reg{FLAGS}, menyebutkan posisi bitnya, dan menjelaskan
          kapan setiap \eng{flag} bernilai 1;
    \item memprediksi perubahan \eng{flag} akibat operasi aritmatika
          dan menjelaskan akibatnya bagi alur program;
    \item menjelaskan mengapa register \reg{CS}, \reg{IP}, \reg{SS},
          dan \reg{SP} tidak boleh diubah sembarangan.
\end{enumerate}
\end{tujuanbox}

\begin{materibox}
\begin{enumerate}
    \item Register umum \reg{AX}, \reg{BX}, \reg{CX}, \reg{DX} dan
          akses byte rendah/tinggi (\reg{AL}/\reg{AH},
          \reg{BL}/\reg{BH}, \reg{CL}/\reg{CH}, \reg{DL}/\reg{DH}).
    \item Register segmen \reg{CS}, \reg{DS}, \reg{SS}, \reg{ES} dan
          perannya pada pengalamatan memori.
    \item Register \eng{pointer} dan \eng{index}: \reg{SP}, \reg{BP},
          \reg{SI}, \reg{DI}, serta register \reg{IP}.
    \item Register \eng{flag} 16-bit: \flag{CF}, \flag{PF},
          \flag{AF}, \flag{ZF}, \flag{SF}, \flag{TF}, \flag{IF},
          \flag{DF}, \flag{OF} beserta posisi bitnya.
    \item Pengaruh \eng{flag} terhadap keputusan program melalui
          lompatan bersyarat.
    \item Register yang tidak boleh diubah sembarangan dan
          akibatnya.
\end{enumerate}
\end{materibox}

\subfile{section-01}
\subfile{section-02}
\subfile{section-03}

\section{Contoh Program}
\label{sec:contoh-bab5}

Bagian ini menyajikan tiga contoh. Dua contoh pertama dapat langsung
dirakit dan dijalankan, sedangkan contoh ketiga adalah potongan
singkat untuk memperlihatkan perilaku register pada tingkat byte.

\subsection{Contoh 1: Memprediksi Flag Sebelum Menjalankan Program}

Program pada Kode~\ref{lst:prediksi-flag} menjalankan dua
penjumlahan yang sengaja dipilih untuk memicu kondisi \eng{flag}
berbeda. Sebelum program dijalankan, mahasiswa diminta memprediksi
nilai \flag{CF}, \flag{OF}, \flag{SF}, dan \flag{ZF} pada
Tabel~\ref{tab:prediksi-flag}, lalu memeriksa hasil prediksinya
melalui \eng{Single Step}.

\lstinputlisting[
    language=MASM,
    caption={Program \file{prediksi-flag.asm} untuk mengamati
    perubahan \eng{flag}.},
    label={lst:prediksi-flag}
]{\CodeDir/prediksi-flag.asm}

\begin{table}[H]
    \centering
    \caption{Lembar prediksi \eng{flag} untuk
    Kode~\ref{lst:prediksi-flag}. Isi kolom terakhir berdasarkan
    pengamatan pada emulator.}
    \label{tab:prediksi-flag}
    \small
    \begin{tabular}{@{}clccccc@{}}
        \toprule
        \textbf{Kasus} & \textbf{Operasi} & \flag{CF} & \flag{OF} &
        \flag{SF} & \flag{ZF} & \textbf{Hasil (\reg{AX})} \\
        \midrule
        1 & \hex{007F} $+$ \hex{0001} & \dots & \dots & \dots &
        \dots & \dots \\
        \addlinespace
        2 & \hex{8000} $+$ \hex{8000} & \dots & \dots & \dots &
        \dots & \dots \\
        \bottomrule
    \end{tabular}
\end{table}

Penjelasan potongan penting pada program tersebut adalah sebagai
berikut.

\begin{itemize}
    \item \lstasm|MOV AX, 007Fh| memuat $+127$ ke akumulator tanpa
          mengubah \eng{flag}.
    \item \lstasm|ADD AX, BX| adalah instruksi yang mengubah
          \eng{flag}. Kedua operand bernilai positif, tetapi
          hasilnya \hex{0080} yang bit tandanya menyala. Inilah
          \eng{overflow} bertanda, sehingga \flag{OF} bernilai 1.
    \item \lstasm|MOV hasil1, AX| menyimpan salinan hasil ke segmen
          data. Karena \asm{MOV} tidak menyentuh \eng{flag}, nilai
          \eng{flag} hasil penjumlahan masih dapat diamati pada
          \eng{debugger} setelah baris ini dijalankan.
    \item Kasus kedua memakai nilai \hex{8000} pada kedua operand.
          Hasilnya nol dengan simpanan keluar, sehingga
          \flag{CF} dan \flag{ZF} bernilai 1 sekaligus.
\end{itemize}

\begin{session}[caption={Keluaran program \file{prediksi-flag.asm}.},
label={lst:keluaran-prediksi-flag}]
Pengamatan flag selesai.
\end{session}

\subsection{Contoh 2: Memilih Pesan Berdasarkan Flag}

Contoh pertama memperlihatkan bagaimana \eng{flag} terbentuk.
Contoh kedua, pada Kode~\ref{lst:pesan-flag}, memperlihatkan
bagaimana \eng{flag} \emph{dipakai}. Program menjumlahkan
\hex{F0} dan \hex{10} pada register \reg{AL} lalu mencetak pesan yang
berbeda sesuai keadaan \eng{flag}.

\lstinputlisting[
    language=MASM,
    caption={Program \file{pesan-flag.asm} memilih pesan berdasarkan
    \eng{flag} hasil \lstasm{ADD AL, BL}.},
    label={lst:pesan-flag}
]{\CodeDir/pesan-flag.asm}

Bagian yang perlu diperhatikan adalah urutan pemeriksaan
\eng{flag}. Setiap kali \lstasm|ADD AL, BL| selesai, register
\reg{FLAGS} memuat seluruh penanda hasil operasi tersebut sekaligus.
Program kemudian memeriksanya satu per satu:

\begin{enumerate}
    \item \lstasm|JNC tanpa_carry| melompat bila \flag{CF} bernilai
          0. Bila \flag{CF} bernilai 1, instruksi berikutnya
          mencetak pesan carry.
    \item \lstasm|JNZ bukan_nol| melompat bila \flag{ZF} bernilai 0.
          Bila hasilnya nol, program mencetak pesan hasil nol.
    \item \lstasm|JNS tanpa_negatif| melompat bila \flag{SF} bernilai
          0. Bila bit tanda menyala, program mencetak pesan hasil
          negatif.
    \item \lstasm|CMP AL, 0| dipakai untuk membangkitkan kembali
          \eng{flag} berdasarkan isi \reg{AL} saat itu. Instruksi
          ini tidak menyimpan hasil pengurangan, hanya meninggalkan
          jejak pada \eng{flag}.
\end{enumerate}

Karena \hex{F0} ditambah \hex{10} menghasilkan \hex{100}, register
\reg{AL} menyimpan \hex{00} dengan \flag{CF} bernilai 1 dan
\flag{ZF} bernilai 1. Keluaran program tampak pada
Sesi~\ref{lst:keluaran-pesan-flag}.

\begin{session}[caption={Keluaran program \file{pesan-flag.asm}.},
label={lst:keluaran-pesan-flag}]
Terjadi carry, CF = 1.
Hasil penjumlahan nol, ZF = 1.
\end{session}

\begin{infobox}
Layanan \code{INT 21h} tidak merusak \eng{flag} hasil penjumlahan.
Instruksi \asm{INT} menyimpan \reg{FLAGS} ke \eng{stack} dan
\asm{IRET} mengembalikannya, sehingga pemeriksaan \flag{CF},
\flag{ZF}, dan \flag{SF} secara berurutan tetap sah.
\end{infobox}

\subsection{Contoh 3: Akses Byte Rendah dan Byte Tinggi}

Potongan pada Kode~\ref{lst:akses-byte} memperlihatkan bahwa satu
register 16 bit dapat diperlakukan sebagai dua register 8 bit yang
mandiri.

\begin{lstlisting}[
    language=MASM,
    caption={Mengakses byte rendah dan byte tinggi register
    \reg{AX}.},
    label={lst:akses-byte}
]
.MODEL SMALL
.STACK 100h

.DATA
    dummy DB 0          ; variabel penanda agar segmen data ada

.CODE
START:
    MOV AX, @DATA
    MOV DS, AX

    MOV AX, 1234h       ; AH = 12h, AL = 34h
    MOV BL, AH          ; BL = 12h (salin byte tinggi)
    MOV BH, AL          ; BH = 34h (salin byte rendah)

    ADD AL, BL          ; AL = 34h + 12h = 46h, AH tidak berubah
    MOV CL, AL          ; CL = 46h

    MOV AH, 4Ch
    MOV AL, 00h
    INT 21h

END START
\end{lstlisting}

Setelah seluruh instruksi di atas dijalankan, isi register adalah
\reg{AX} = \hex{1246}, \reg{BX} = \hex{3412}, dan \reg{CL} =
\hex{46}. Perhatikan bahwa \lstasm|ADD AL, BL| hanya mengubah
\reg{AL}; nilai \reg{AH} tetap \hex{12} karena \eng{flag} dan
\reg{AH} tidak tersentuh oleh operasi 8 bit pada \reg{AL}.

\section{Praktikum}
\label{sec:praktikum-bab5}

\begin{praktikumbox}
\textbf{Judul:} Pengamatan register dan \eng{flag} dengan
\eng{Single Step} pada Emu8086.\\[0.3em]
\textbf{Sub-CPMK:} 1.2 \quad\textbullet\quad
\textbf{Alokasi waktu:} 100 menit \quad\textbullet\quad
\textbf{Bentuk:} individu
\end{praktikumbox}

\subsection{Tujuan Praktikum}

Setelah menyelesaikan praktikum ini, mahasiswa mampu mengamati isi
setiap register 8086, membaca keadaan \eng{flag} setelah operasi
aritmatika, dan menjelaskan akibat perubahan nilai register secara
manual terhadap jalannya program.

\subsection{Alat dan Bahan}

\begin{itemize}
    \item Komputer dengan Emu8086 terpasang.
    \item Berkas \file{prediksi-flag.asm} dan \file{pesan-flag.asm}
          dari direktori \file{code/asm/}.
    \item Lembar kerja pengamatan register dan \eng{flag}.
\end{itemize}

\subsection{Langkah Kerja}

\begin{enumerate}
    \item \textbf{Menyiapkan program pengamatan.}
    \begin{enumerate}
        \item Buka Emu8086, lalu buka berkas
              \file{prediksi-flag.asm}.
        \item Tekan \file{Compile and Emulate}, kemudian tahan pada
              jendela emulator tanpa menjalankan program.
        \item Pastikan jendela \file{Registers} dan jendela
              \file{Flags} terlihat. Bila belum, aktifkan melalui
              menu \file{View}.
    \end{enumerate}

    \item \textbf{Mengamati register dengan \eng{Single Step}.}
    \begin{enumerate}
        \item Tekan tombol \file{Single Step} satu kali untuk
              menjalankan instruksi pertama.
        \item Catat nilai \reg{AX}, \reg{BX}, \reg{CX}, \reg{DX},
              \reg{DS}, \reg{SP}, dan \reg{IP} pada
              Tabel~\ref{tab:pengamatan-register} setiap kali satu
              instruksi selesai.
        \item Berhenti pada instruksi \lstasm|ADD AX, BX| dan catat
              keadaan \eng{flag} sebelum instruksi dijalankan.
        \item Jalankan instruksi \lstasm|ADD AX, BX| satu langkah,
              lalu catat keadaan \eng{flag} sesudahnya.
        \item Lanjutkan sampai program selesai dan catat keadaan
              \eng{flag} pada akhir kasus pertama dan kasus kedua.
    \end{enumerate}

    \item \textbf{Mengisi tabel pengamatan \eng{flag}.}
    \begin{enumerate}
        \item Isi kolom prediksi pada Tabel~\ref{tab:prediksi-flag}
              \emph{sebelum} menjalankan kasus yang bersangkutan.
        \item Bandingkan prediksi dengan hasil pengamatan.
        \item Untuk setiap prediksi yang salah, tuliskan penjelasan
              mengapa \eng{flag} tersebut bernilai demikian.
    \end{enumerate}

    \item \textbf{Mengubah nilai register secara manual.}
    \begin{enumerate}
        \item Pada jendela \file{Registers}, klik ganda nilai
              \reg{AX}, ubah menjadi \hex{FFFF}, lalu tekan
              \file{Enter}.
        \item Jalankan satu instruksi berikutnya dan amati pengaruh
              nilai baru tersebut terhadap hasil dan \eng{flag}.
        \item Ubah nilai \reg{AL} saja menjadi \hex{00} dan amati
              bahwa \reg{AH} tidak ikut berubah.
        \item Ulangi langkah serupa untuk \reg{BX} dan catat
              perbedaannya.
    \end{enumerate}

    \item \textbf{Menguji pengaruh perubahan register kendali.}
    \begin{enumerate}
        \item Jalankan berkas \file{pesan-flag.asm} dengan
              \eng{Single Step}.
        \item Pada jendela \file{Flags}, ubah nilai \flag{CF} secara
              manual dari 1 menjadi 0, lalu lanjutkan eksekusi.
        \item Amati pesan mana yang tercetak dan jelaskan sebabnya
              berdasarkan urutan pemeriksaan \eng{flag}.
        \item Kembalikan nilai \flag{CF} dan ulangi percobaan untuk
              \flag{ZF}.
    \end{enumerate}

    \item \textbf{Mengamati register \reg{SP} dan \reg{SS}.}
    \begin{enumerate}
        \item Tuliskan potongan program singkat berisi
              \lstasm|MOV AX, 1234h| lalu \lstasm|PUSH AX|.
        \item Jalankan program dengan \eng{Single Step} dan catat
              nilai \reg{SP} sebelum dan sesudah \asm{PUSH}.
        \item Jelaskan mengapa \reg{SP} berubah sebesar 2, bukan 1.
    \end{enumerate}
\end{enumerate}

\subsection{Tabel Pengamatan}

\begin{table}[H]
    \centering
    \caption{Lembar pengamatan isi register pada praktikum bab ini.}
    \label{tab:pengamatan-register}
    \small
    \begin{tabular}{@{}lp{2.0cm}p{2.0cm}p{2.0cm}p{2.0cm}@{}}
        \toprule
        \textbf{Instruksi} & \reg{AX} & \reg{BX} & \reg{SP} &
        \reg{IP} \\
        \midrule
        \lstasm|MOV AX, @DATA| & \dots & \dots & \dots & \dots \\
        \addlinespace
        \lstasm|MOV DS, AX| & \dots & \dots & \dots & \dots \\
        \addlinespace
        \lstasm|MOV AX, 007Fh| & \dots & \dots & \dots & \dots \\
        \addlinespace
        \lstasm|MOV BX, 0001h| & \dots & \dots & \dots & \dots \\
        \addlinespace
        \lstasm|ADD AX, BX| & \dots & \dots & \dots & \dots \\
        \addlinespace
        \lstasm|MOV hasil1, AX| & \dots & \dots & \dots & \dots \\
        \bottomrule
    \end{tabular}
\end{table}

\subsection{Hasil yang Dikumpulkan}

\begin{enumerate}
    \item Tabel~\ref{tab:pengamatan-register} yang telah terisi
          lengkap.
    \item Tabel~\ref{tab:prediksi-flag} yang telah terisi prediksi
          dan hasil pengamatan.
    \item Tangkapan layar jendela \file{Registers} dan
          \file{Flags} pada saat \lstasm|ADD AX, BX| selesai
          dijalankan.
    \item Catatan hasil percobaan mengubah \flag{CF} dan \flag{ZF}
          secara manual beserta penjelasan sebabnya.
\end{enumerate}

\begin{asesmenbox}
\textbf{Asesmen praktikum 3 (bobot 5\% dari komponen praktikum).}

\begin{tabular}{@{}p{6.2cm}cp{4.6cm}@{}}
    \toprule
    \textbf{Indikator} & \textbf{Bobot} & \textbf{Kriteria} \\
    \midrule
    Ketepatan pengamatan register & 30\% &
    Seluruh nilai register tercatat benar pada setiap langkah \\
    \addlinespace
    Ketepatan pengamatan \eng{flag} & 30\% &
    Nilai \flag{CF}, \flag{OF}, \flag{SF}, \flag{ZF} sesuai hasil
    pengamatan \\
    \addlinespace
    Analisis perubahan manual & 25\% &
    Penjelasan sebab perubahan hasil dan \eng{flag} logis dan tepat \\
    \addlinespace
    Kerapian dokumentasi & 15\% &
    Lembar kerja lengkap, rapi, dan dapat diulang \\
    \bottomrule
\end{tabular}
\end{asesmenbox}

\section{Latihan}
\label{sec:latihan-bab5}

\begin{enumerate}
    \item \textbf{Identifikasi.} Sebutkan keempat belas register
          8086 dan kelompokkan menjadi register umum, segmen,
          \eng{pointer}/\eng{index}, kendali, dan status.

    \item \textbf{Akses byte.} Register \reg{AX} semula berisi
          \hex{1234}. Tentukan isi \reg{AH} dan \reg{AL}. Setelah
          \lstasm|MOV AL, 20h| dijalankan, tentukan isi \reg{AX}.
          Jelaskan mengapa \reg{AH} tidak berubah.

    \item \textbf{Peran khas.} Jelaskan peran khas \reg{AX},
          \reg{BX}, \reg{CX}, dan \reg{DX}. Mengapa peran itu ada
          walaupun keempat register sama-sama berukuran 16 bit?

    \item \textbf{Register segmen.} Jelaskan peran \reg{CS},
          \reg{DS}, \reg{SS}, dan \reg{ES}. Tuliskan pasangan alamat
          yang dibentuk oleh masing-masing register segmen.

    \item \textbf{Pointer dan index.} Jelaskan tugas \reg{SP},
          \reg{BP}, \reg{SI}, \reg{DI}, dan \reg{IP}. Sebutkan
          pasangan segmen untuk setiap register tersebut.

    \item \textbf{Flag.} Jelaskan arti \flag{CF}, \flag{PF},
          \flag{AF}, \flag{ZF}, \flag{SF}, \flag{OF}, \flag{TF},
          \flag{IF}, dan \flag{DF}. Kelompokkan menjadi \eng{flag}
          status dan \eng{flag} kendali.

    \item \textbf{Prediksi.} Diberikan dua kasus penjumlahan 16 bit.
          \begin{enumerate}
              \item \hex{7FFF} $+$ \hex{0001}
              \item \hex{8000} $+$ \hex{8000}
          \end{enumerate}
          Untuk setiap kasus, tentukan hasil 16 bit serta nilai
          \flag{CF}, \flag{OF}, \flag{SF}, dan \flag{ZF}. Jelaskan
          alasan setiap nilai.

    \item \textbf{Keputusan.} Setelah \lstasm|CMP AL, 0|, program
          menjalankan \lstasm|JNZ bukan_nol|. Jelaskan \eng{flag}
          yang diperiksa, kapan lompatan terjadi, dan apa artinya
          bagi alur program.
\end{enumerate}

\begin{refleksibox}
\vspace{1.1em}
Jawab pertanyaan berikut pada jurnal praktikum Anda.
\begin{enumerate}
    \item Bagian mana dari pembagian register yang paling sulit Anda
          ingat? Cara apa yang Anda pakai untuk mengingatnya?
    \item Setelah mengamati \eng{flag} dengan \eng{single step},
          apakah Anda dapat memprediksi hasil operasi berikutnya
          dengan lebih baik? Jelaskan.
    \item Konsep apa yang perlu Anda ulangi sebelum mempelajari
          segmentasi memori pada bab berikutnya?
\end{enumerate}
\end{refleksibox}

\section{Rangkuman}
\label{sec:rangkuman-bab5}

Pada bab ini telah dibahas:

\begin{itemize}
    \item Register adalah tempat penyimpanan tercepat di dalam
          prosesor. Prosesor 8086 memiliki empat belas register
          16 bit yang dapat dipakai programmer.

    \item Register umum \reg{AX}, \reg{BX}, \reg{CX}, dan \reg{DX}
          dapat diakses sebagai dua byte terpisah: \reg{AH}/\reg{AL},
          \reg{BH}/\reg{BL}, \reg{CH}/\reg{CL}, dan \reg{DH}/\reg{DL}.
          Menulis ke byte rendah tidak mengubah byte tinggi.

    \item Setiap register umum memiliki peran khas sebagai akibat
          cara penyandian instruksi: \reg{AX} sebagai akumulator,
          \reg{BX} sebagai basis pengalamatan, \reg{CX} sebagai
          pencacah, dan \reg{DX} sebagai pasangan data pada operasi
          32 bit serta penunjuk alamat port.

    \item Register segmen \reg{CS}, \reg{DS}, \reg{SS}, dan \reg{ES}
          menyimpan nomor segmen, bukan data. Pasangannya adalah
          \reg{IP} untuk kode, \reg{BP}/\reg{SP} untuk \eng{stack},
          serta \reg{SI} dan \reg{DI} untuk data dan tujuan operasi
          string.

    \item Register \eng{pointer} dan \eng{index} menyimpan
          \eng{offset}. \reg{SP} menunjuk puncak \eng{stack} dan
          berubah otomatis pada \asm{PUSH}/\asm{POP}, \reg{BP}
          menjadi acuan tetap di dalam \eng{stack}, sedangkan
          \reg{SI} dan \reg{DI} menunjuk sumber dan tujuan operasi
          string.

    \item Register \reg{IP} menyimpan \eng{offset} instruksi
          berikutnya dan selalu berpasangan dengan \reg{CS}.

    \item Register \reg{FLAGS} 16 bit memuat sembilan \eng{flag}
          aktif: \flag{CF}, \flag{PF}, \flag{AF}, \flag{ZF},
          \flag{SF}, \flag{OF} sebagai \eng{flag} status, serta
          \flag{TF}, \flag{IF}, \flag{DF} sebagai \eng{flag}
          kendali.

    \item \flag{CF} melaporkan simpanan keluar pada bilangan tak
          bertanda, sedangkan \flag{OF} melaporkan \eng{overflow}
          pada bilangan bertanda. Keduanya harus dibedakan dengan
          cermat.

    \item \eng{Flag} menjadi berguna ketika dibaca oleh lompatan
          bersyarat, sedangkan \asm{MOV}, \asm{LEA}, \asm{PUSH}, dan
          \asm{POP} tidak mengubah \eng{flag}.

    \item Register \reg{CS}, \reg{IP}, \reg{SS}, dan \reg{SP}
          mengendalikan alur program dan letak \eng{stack}, sehingga
          tidak boleh diubah sembarangan.
\end{itemize}

\begin{checklistbox}
Tandai setiap butir yang sudah Anda kuasai.

\begin{itemize}
    \item[\checkbox] Saya dapat menyebutkan keempat belas register
          8086 beserta kelompoknya.
    \item[\checkbox] Saya dapat menjelaskan akses byte rendah dan
          byte tinggi pada \reg{AX}, \reg{BX}, \reg{CX}, \reg{DX}.
    \item[\checkbox] Saya dapat menjelaskan peran khas setiap
          register umum beserta alasannya.
    \item[\checkbox] Saya dapat menjelaskan peran \reg{CS}, \reg{DS},
          \reg{SS}, dan \reg{ES} dalam pengalamatan memori.
    \item[\checkbox] Saya dapat menjelaskan perbedaan tugas \reg{SP},
          \reg{BP}, \reg{SI}, \reg{DI}, dan \reg{IP}.
    \item[\checkbox] Saya dapat menyebutkan sembilan \eng{flag}
          beserta posisi bitnya.
    \item[\checkbox] Saya dapat memprediksi \flag{CF}, \flag{OF},
          \flag{SF}, dan \flag{ZF} pada operasi \asm{ADD}.
    \item[\checkbox] Saya dapat menjelaskan akibat mengubah
          \reg{CS}, \reg{IP}, \reg{SS}, atau \reg{SP} secara
          sembarangan.
\end{itemize}

\vspace{0.4em}
Bila masih ada butir yang belum bertanda, ulangi bagian yang
bersesuaian sebelum melanjutkan ke bab berikutnya.
\end{checklistbox}

\begin{infobox}
\textbf{Rujukan bab.} Uraian mengenai fungsi dan penyandian register
8086 bersumber dari \cite{intel1979}. Penjelasan tentang \eng{flag}
dan pengaruhnya terhadap eksekusi diperkuat oleh \cite{brey2009}
dan \cite{hyde2010}. Ringkasan praktis mengenai daftar register dan
\eng{flag} dapat ditelusuri pada \cite{wikibooksx86}.
\end{infobox}

\section{Evaluasi}
\label{sec:evaluasi-bab5}

\subsection{Soal Pilihan Ganda}

\begin{enumerate}
    \item Setelah \lstasm|MOV AL, 20h| dijalankan pada register
          \reg{AX} yang semula berisi \hex{1234}, isi \reg{AX}
          menjadi \dots
          \begin{enumerate}
              \item \hex{1220} \item \hex{2034}
              \item \hex{0020} \item \hex{2020}
          \end{enumerate}

    \item Register yang berpasangan dengan \reg{CS} untuk menentukan
          alamat instruksi berikutnya adalah \dots
          \begin{enumerate}
              \item \reg{SP} \item \reg{IP}
              \item \reg{DS} \item \reg{DI}
          \end{enumerate}

    \item Instruksi \lstasm|ADD AL, BL| dengan \reg{AL} = \hex{7F}
          dan \reg{BL} = \hex{01} menghasilkan \dots
          \begin{enumerate}
              \item \flag{CF} = 1 dan \flag{OF} = 0
              \item \flag{CF} = 0 dan \flag{OF} = 1
              \item \flag{CF} = 1 dan \flag{OF} = 1
              \item \flag{CF} = 0 dan \flag{OF} = 0
          \end{enumerate}

    \item Register yang menunjuk puncak \eng{stack} adalah \dots
          \begin{enumerate}
              \item \reg{BP} \item \reg{SP}
              \item \reg{SI} \item \reg{DI}
          \end{enumerate}

    \item \eng{Flag} yang bernilai 1 ketika hasil operasi bernilai
          nol adalah \dots
          \begin{enumerate}
              \item \flag{CF} \item \flag{SF}
              \item \flag{ZF} \item \flag{OF}
          \end{enumerate}
\end{enumerate}

\subsection{Soal Uraian}

\begin{enumerate}
    \item Jelaskan mengapa \lstasm|MOV AL, 5| dan \lstasm|MOV AX, 5|
          menghasilkan keadaan register yang berbeda bila sebelumnya
          \reg{AH} tidak bernilai nol. Sertakan contoh angka.

    \item Sebutkan peran khas \reg{BX}, \reg{CX}, dan \reg{DX}.
          Jelaskan mengapa peran tersebut muncul walaupun keempat
          register umum sama-sama dapat menampung nilai.

    \item Diberikan keadaan awal \reg{AL} = \hex{80} dan
          \reg{BL} = \hex{80}. Tentukan hasil
          \lstasm|ADD AL, BL| beserta nilai \flag{CF}, \flag{OF},
          \flag{SF}, dan \flag{ZF}. Jelaskan mengapa \flag{CF} dan
          \flag{OF} bernilai demikian.

    \item Jelaskan peran \reg{CS}, \reg{DS}, \reg{SS}, dan \reg{ES}
          beserta pasangan alamatnya. Jelaskan pula akibat yang
          mungkin terjadi bila \reg{SS} dan \reg{SP} diubah secara
          sembarangan.
\end{enumerate}

\subsection{Kisi-kisi Penilaian}

\begin{table}[H]
    \centering
    \caption{Kisi-kisi evaluasi Bab~\ref{chap:register}.}
    \label{tab:kisi-bab5}
    \small
    \begin{tabular}{@{}p{1.5cm}p{5.0cm}p{2.6cm}c@{}}
        \toprule
        \textbf{Sub-CPMK} & \textbf{Indikator yang diukur} &
        \textbf{Bentuk soal} & \textbf{Bobot} \\
        \midrule
        1.3 & Mengidentifikasi register umum dan akses byte &
        Pilihan ganda 1; Uraian 1, 2 & 25\% \\
        \addlinespace
        1.3 & Menjelaskan register segmen dan pasangannya &
        Pilihan ganda 2; Uraian 4 & 25\% \\
        \addlinespace
        1.3 & Menjelaskan register \eng{pointer} dan \eng{index} &
        Pilihan ganda 4 & 15\% \\
        \addlinespace
        1.3 & Menganalisis perubahan \eng{flag} hasil operasi &
        Pilihan ganda 3, 5; Uraian 3 & 35\% \\
        \bottomrule
    \end{tabular}
\end{table}

\begin{checklistbox}
\vspace{1.1em}
\textbf{Checklist Pencapaian Kompetensi Bab~\ref{chap:register}}

Seorang mahasiswa dinyatakan kompeten pada Sub-CPMK 1.3 apabila:
\begin{itemize}
    \item[\checked] dapat menyebutkan register umum beserta byte
          tinggi dan byte rendahnya;
    \item[\checked] dapat menjelaskan peran khas AX, BX, CX, dan DX;
    \item[\checked] dapat menjelaskan peran CS, DS, SS, dan ES pada
          pengalamatan memori;
    \item[\checked] dapat menjelaskan tugas SP, BP, SI, DI, dan IP;
    \item[\checked] dapat menyebutkan sembilan \eng{flag} dan kapan
          masing-masing bernilai 1;
    \item[\checked] dapat memprediksi perubahan \eng{flag} pada
          operasi penjumlahan;
    \item[\checked] dapat menjelaskan pengaruh \eng{flag} terhadap
          keputusan program.
\end{itemize}
\end{checklistbox}


\end{document}
